% !PRINTED 108
\weberpage{108}
Beginnen wir also mit den cubischen Gleichungen, und verstehen
 unter $x_0, x_1, x_2$ die Wurzeln irgend einer \gesperrt{Abel'schen}
Gleichung dritten Grades. Da 3 eine Primzahl ist, so muss die
Gleichung cyklisch sein, d.\,h. \gesperrt{die cyklischen Functionen
der Grössen} $x_0, x_1, x_2$ \gesperrt{sind rationale Zahlen.}

Bezeichnen wir mit $\varrho$ eine imaginäre dritte Einheitswurzel,
so haben wir die beiden Resolventen
\weq{1}{\begin{aligned}
(\varrho, x_0) &= x_0 + \varrho\, x_1 + \varrho^2 x_2\\
(\varrho^2, x_0) &= x_0 + \varrho^2 x_1 + \varrho\, x_2,
\end{aligned}}
deren Cuben also. Zahlen des Körpers $R(\varrho)$ und deren Product % [sic: Weber prints a full stop after ``also'', although the sentence continues.]
eine rationale Zahl sein muss.

Wir setzen demnach, indem wir mit $a, b, c$ rationale (ganze
oder gebrochene) Zahlen bezeichnen,
\weq{2}{\begin{aligned}
(\varrho, x_0)^3 &= a + b\varrho\\
(\varrho^2, x_0)^3 &= a + b\varrho^2\\
(\varrho, x_0)(\varrho^2, x_0) &= c,
\end{aligned}}
also auch
\weq{3}{(a+b\varrho)(a+b\varrho^2) = c^3.}

Nun haben wir im Körper $R(\varrho)$ ausser den sechs Einheiten
\[ \pm 1,\ \pm\varrho,\ \pm\varrho^2, \]
die alle als Potenzen von $-\varrho$ dargestellt werden können, die
Primzahlen $\sqrt{-3} = \varrho - \varrho^2$, die reellen Primzahlen $q$ von der
Form $3N+2$ und die beiden complexen Factoren der reellen
Primzahlen $p$ von der Form $3N+1$. Die Zerlegung dieser
Primzahlen in die beiden complexen Factoren $\pi, \pi'$:
\weq{4}{p = \pi\pi'}
haben wir im \S.\,172 des ersten Bandes dargestellt in der Form
\weq{5}{p = \psi_1(\varrho)\,\psi_1(\varrho^2).}
Wir können also
\weq{6}{\pi = \psi_1(\varrho), \quad \pi' = \psi_1(\varrho^2),}
setzen und haben dadurch unter den verschiedenen zu $\pi$ associirten
 Zahlen eine bestimmte ausgewählt, und diese Zahlen
$\pi, \pi'$ stehen in einer bestimmten Beziehung zu den Kreis\-theilungsperioden
 der $p^{\text{ten}}$ Einheitswurzeln $\eta$:
\weq{7}{(\varrho, \eta)^3 = p\pi, \quad (\varrho^2, \eta)^3 = p\pi'}
[Bd.\ I, \S.\,172, (5)].
